\documentclass[11pt,a4paper]{article}

\usepackage[T1]{fontenc}
\usepackage[utf8]{inputenc}
\usepackage{lmodern}
\usepackage{microtype}
\usepackage[a4paper,margin=25mm,headheight=14pt]{geometry}
\usepackage{setspace}
\usepackage{enumitem}
\usepackage{booktabs}
\usepackage{tabularx}
\usepackage{longtable}
\usepackage{array}
\usepackage{xcolor}
\usepackage{amsmath,amssymb}
\usepackage{fancyhdr}
\usepackage{xurl}
\usepackage{hyperref}

\definecolor{CQGblue}{RGB}{31,74,125}
\definecolor{DarkGreen}{RGB}{35,105,65}
\definecolor{DarkRed}{RGB}{145,38,38}
\definecolor{Amber}{RGB}{160,94,0}
\definecolor{LightBlue}{RGB}{235,242,250}
\definecolor{LightGreen}{RGB}{236,247,239}
\definecolor{LightAmber}{RGB}{252,245,229}

\hypersetup{
  colorlinks=true,
  linkcolor=CQGblue,
  urlcolor=CQGblue,
  pdftitle={Updated editorial assessment of Paper I for Classical and Quantum Gravity},
  pdfauthor={Codex editorial review},
  pdfsubject={Pre-submission editorial assessment},
  pdfkeywords={Classical and Quantum Gravity; editorial review; controlled rail; Vaidya; brachistochrone}
}

\setstretch{1.05}
\setlength{\parindent}{0pt}
\setlength{\parskip}{0.47em}
\setlist[itemize]{leftmargin=1.65em,itemsep=0.32em,topsep=0.35em}
\setlist[enumerate]{leftmargin=1.9em,itemsep=0.38em,topsep=0.4em}
\renewcommand{\arraystretch}{1.17}

\pagestyle{fancy}
\fancyhf{}
\lhead{Updated editorial review --- Paper I}
\rhead{CQG pre-submission assessment}
\cfoot{\thepage}

\newcommand{\must}{\textcolor{DarkRed}{\textbf{Must fix}}}
\newcommand{\should}{\textcolor{Amber}{\textbf{Should fix}}}
\newcommand{\preserve}{\textcolor{DarkGreen}{\textbf{Preserve}}}
\newcommand{\paper}{\emph{Controlled-rail brachistochrones in non-stationary spacetimes: conformal symmetry, Kodama energy, and Vaidya dynamics}}
\newcommand{\Ehat}{\hat E}
\newcommand{\SD}{S_{\mathcal D}}

\begin{document}
\hypersetup{pageanchor=false}

\begin{titlepage}
  \centering
  \vspace*{1.5cm}
  {\Large\bfseries Updated editorial assessment for\\[0.35em]
  \emph{Classical and Quantum Gravity}\par}
  \vspace{1.2cm}
  {\LARGE\bfseries \paper\par}
  \vspace{0.8cm}
  {\large Complete review of the final 34-page production build\par}
  \vspace{1.3cm}
  \begin{minipage}{0.87\textwidth}
    \colorbox{LightAmber}{\parbox{0.93\textwidth}{
      \textbf{Simulated decision: short mandatory editorial revision, then external review.}
      The manuscript is now a credible CQG Research Paper and the earlier major editorial
      blockers have largely been removed. It should not be uploaded unchanged, however:
      a compact set of notation, self-containedness, bibliography and archive-version
      inconsistencies remains. None requires reopening the central off-shell derivation.
    }}
  \end{minipage}
  \vfill
  {\large Prepared as an Associate-Editor-style assessment, not as a substitute for\\
  specialist mathematical peer review\par}
  \vspace{0.45cm}
  {\large 4 August 2026\par}
  \vspace{1.1cm}
  {\small Codex-origin report. The manuscript and its source were not edited during this assessment.\par}
\end{titlepage}

\hypersetup{pageanchor=true}
\pagenumbering{roman}
\tableofcontents
\clearpage
\pagenumbering{arabic}

\section{Mandate, material reviewed, and limits}

This report follows the supplied ``Instructions for Codex -- CQG Editorial Review''. The
mandate is to read the manuscript completely as a simulated Associate Editor, assess its
editorial readiness for \emph{Classical and Quantum Gravity} (CQG), and produce a full
editorial report rather than a new mathematical referee report. The assessment therefore
concentrates on scope, novelty presentation, readability, claim calibration, structure,
figures, notation, bibliography, reproducibility and submission compliance.

The authoritative manuscript is the current file \path{paper1/paper1.pdf}, compiled on
4 August 2026. It contains 34 pages, nine numbered main figures, two appendix figures,
two tables and 43 references. I read all pages, including Appendices A--C, the data
availability statement, acknowledgements, funding and competing-interest statements,
ORCID line and complete bibliography. I also inspected the current \LaTeX{} source,
the production log, the result-to-script table, the cover-letter source, the prior
editorial review and handoff documents, the public Zenodo/DataCite metadata for the three
DOIs used in the submission package, and the current official CQG/IOP author guidance.

The complete PDF was rendered and visually inspected page by page. Layout findings in this
report therefore refer to the actual final rendering, not merely to extracted text. The
build has no undefined references or citations, no duplicate labels and no overfull boxes.

This remains an editorial assessment. I do not independently certify every Hamiltonian
sign, the proof of normality, the exact residue formulae, the genus-two period
construction, or the asserted error bounds. Those questions require specialist referees.
Where the manuscript makes mutually inconsistent statements, uses an undefined object,
or claims that a formula is printed when it is not, that is an editorially verifiable
defect and is treated as such.

\subsection*{Severity vocabulary}

\begin{description}[leftmargin=2.1cm,style=nextline]
  \item[\must] A concrete inconsistency likely to reduce referee confidence or obstruct
  reproducibility. It should be corrected before submission.
  \item[\should] A readability, positioning or production improvement that is unlikely to
  determine validity by itself but materially improves desk-review prospects.
  \item[\preserve] A feature already working well that should not be lost during cleanup.
\end{description}

\clearpage
\section{Executive summary}

The manuscript asks a clear question of genuine interest to CQG: how should a relativistic
brachistochrone be formulated when the background geometry evolves and there is no timelike
Killing energy? Its answer is a controlled rail on which the invariant
$-u\!\cdot\!W=\Ehat$ is actively maintained, with $W$ selected through a Killing,
conformal-Killing and, in spherical symmetry, Kodama hierarchy. The paper combines a
Pontryagin formulation, regular-domain existence and normality, a conditional
Hamilton--Jacobi verification criterion, an FLRW base case and a detailed ingoing-Vaidya
application. The most important technical deliverable is the complete first-order
adiabatic response, including the off-shell costate displacement, with a numerical
$O(\varepsilon^2)$ residual test against the true non-autonomous flow.

This is now recognisably a publishable CQG-shaped project. It lies within classical general
relativity, dynamical black holes, constrained Hamiltonian systems and mathematical physics.
It has a memorable conceptual mechanism, physically interpretable Vaidya phenomena,
substantial analytic development and unusually explicit reproducibility support. The
current abstract and conclusion are also disciplined about the difference between proved,
conditional/numerical and open statements.

The earlier review recommended major editorial revision before external review. That
recommendation is no longer appropriate. The current build has addressed most of the
previous blockers:

\begin{itemize}
  \item direct relativistic-brachistochrone and Fermat-principle predecessors are now cited;
  \item Theorem I.5 is stated in the main text before Lemma I.6;
  \item the complete on-shell plus off-shell correction is explicitly distinguished;
  \item Figure 9 reports the fit interval, five samples, residual norm and fitted slopes;
  \item $J_{\rm pen}$ and $J_{\rm deg}$ are mostly separated;
  \item freezing, turning and spectral degeneration are no longer identified;
  \item the Thakurta--Kerr panel has been removed from Figure 1 and preserved separately;
  \item Table A2 no longer overlaps;
  \item all 43 bibliography entries are cited and a References heading is present;
  \item the conclusion is uninterrupted and the page count has fallen from 37 to 34;
  \item PDF metadata is complete and all fonts are embedded.
\end{itemize}

The remaining defects are narrower, but several sit at high-leverage locations: the Kodama
invariant begins from an undefined $T$; the notation changes from $\Ehat$ to $E$ without
definition; $\varepsilon$ is both $\dot m$ and $M\dot m/m$; the off-shell source is both
$\SD$ and $S$ even though $S$ is already the spectral sextic; and Appendix C claims that
a nome series is printed in section 4.4 when it is absent. The public companion-paper title
does not match reference [12], and the cover letter calls an older Zenodo PDF the present
paper. These are not cosmetic in the dismissive sense: they affect whether a referee can
unambiguously reconstruct what has been proved and which archived object is being reviewed.

My simulated editorial decision is therefore:

\begin{center}
\colorbox{LightGreen}{\parbox{0.90\textwidth}{
\textbf{Invite a short author revision, then send the paper to external review.}
Do not desk-reject for scope, length or lack of substance. Do not upload the present files
unchanged. Correct the ten mandatory consistency and snapshot items in section~\ref{sec:must},
rebuild once, and then submit as a CQG Research Paper.
}}
\end{center}

\clearpage
\section{CQG fit and desk-review risk}

CQG states that it welcomes original research across gravitational physics and spacetime
theory, including classical general relativity, black holes, constrained Hamiltonian
systems and differential geometry relevant to gravitation. It also requires Research
Papers to report high-quality original research whose conclusions represent a significant
advance, and asks that articles be interesting to its broad readership and place the result
in the wider context of gravitational physics. The manuscript satisfies the subject-matter
test without strain. Its desk-review risk comes from communication and internal consistency,
not from journal mismatch.

\begin{tabularx}{\textwidth}{@{}p{0.24\textwidth}X p{0.18\textwidth}@{}}
\toprule
Criterion & Current assessment & Risk before fixes \\
\midrule
Scope & Strong fit: dynamical black holes, optical geometry, constrained Hamiltonian
control and exact methods in classical GR. & Low \\
Originality & Distinctive controlled non-autonomous charge, Kodama selector and complete
first-order Vaidya response. External referees must validate priority and mathematics. & Moderate \\
Broad-reader interest & Clear opening question and physical phenomena, but the special-function
layer remains dense and Paper-II references are frequent. & Moderate \\
Claim hierarchy & Much improved; proved, conditional, numerical and open statements are usually
distinguished. Several notation and reconstruction inconsistencies remain. & Moderate \\
Reproducibility & Strong code DOI and script map, but final paper/figure-production files postdate
the cited snapshots. & Moderate \\
Presentation & Clean 34-page PDF, readable figures, no collisions. Visual style is still mixed
between vector publication figures and raster research plots. & Low--moderate \\
Bibliography & Relevant 43-entry list with no padding. Companion title and visible metadata need
correction; DOI presentation is incomplete. & Moderate \\
Ethics/back matter & Data, AI, funding, conflicts and ORCID statements are present and substantively
aligned with current IOP policy. & Low \\
\bottomrule
\end{tabularx}

The length is not itself disqualifying. CQG does not publish a simple page ceiling for Research
Papers in the guidance consulted. The editorial question is whether 34 pages deliver enough
physical and conceptual value for a broad readership. Here the answer can be yes, but only if
the paper stops asking the reader to resolve avoidable symbol collisions and archive mismatches.

The strongest desk-review pitch is not ``a new exact solution'' or ``a new special function''.
It is: a well-defined controlled protocol for non-stationary relativistic travel; a geometrically
motivated selector in spherical symmetry; a concrete Vaidya application; and a complete
first-order response whose need for an off-shell costate term is exposed by a clean convergence
test. The cover letter and abstract should keep that hierarchy.

\clearpage
\section{Principal strengths to preserve}

\subsection{A memorable conceptual spine}

The Killing $\rightarrow$ conformal-Killing $\rightarrow$ Kodama selector hierarchy gives the
paper an identity. It connects stationary optical geometry to dynamical spherical spacetimes
without pretending that the controlled energy is a Noether charge. The hierarchy is strongest
when explicitly restricted: Kodama is the canonical spherical replacement, while in a general
non-spherical spacetime the selector becomes a declared modelling choice.

\subsection{Operational honesty about the rail}

The paper explains that the rail supplies proper acceleration, may do work relative to the
selector, preserves rest mass through $a\!\cdot\!u=0$, and does not minimise fuel or thrust.
This prevents the common objection that arbitrary acceleration would trivialise the problem.
The fixed-charge indicatrix, not unrestricted actuation, defines the admissible velocities.

\subsection{Necessary versus sufficient optimality}

Section 2.2 is editorially mature. It distinguishes Pontryagin extremality from a global HJB
certificate, states regular-domain assumptions, identifies conjugate and Maxwell obstructions,
and admits that no generic global value function is constructed. Figure 2 and Figure A2 now
describe their numerical tests as local perturbation checks. This qualification should remain
untouched.

\subsection{Controlled versus mechanical charges}

The distinction between $J=p_\varphi$ as a reduced-control costate and
$L_{\rm mech}=g(u,\partial_\varphi)$ as mechanical angular momentum is important and clearly
explained. Likewise, the paper correctly says that the ordinary non-autonomous Hamiltonian is
not an interior first integral, while the extended Hamiltonian is the conserved object.

\subsection{The complete first-order correction}

Promoting Theorem I.5 to the main text was the right decision. The momentum-shell shift and the
terminally anchored source make clear why the frozen-family derivative alone leaves an
$O(\varepsilon)$ defect. Figure 9 now functions as a real validation figure: it identifies the
norm, radial interval, five rates, regression variables and slopes 1.003152 and 1.999997. The
unchanged numerical source has a recorded hash. This material is the paper's best response to a
specialist sceptic and should not be compressed out of existence.

\subsection{Physical caveat discipline}

The manuscript consistently states that ingoing Vaidya with $m'<0$ is a formal continuation,
not physical evaporation, and that the full outgoing-Vaidya boundary-value problem remains
open. The conclusion's ``Proved / Conditional and numerical / Outlook'' structure is unusually
helpful. Preserve it almost verbatim in spirit.

\subsection{Reproducibility and disclosure}

Table A2, the Zenodo software release, explicit residuals, script hash, software citations,
data-availability statement, AI disclosure, funding statement, competing-interest declaration
and ORCID collectively exceed the transparency of many theory submissions. The task is now to
make the public snapshots match the final paper, not to reduce the disclosure.

\clearpage
\section{Mandatory corrections before submission}\label{sec:must}

\subsection{Undefined $T$ at the start of the Kodama result}

\must. On PDF page 11, section 4.1 starts from
\[
  -g(T,T)=(f u^v-1)^2/\Ehat^2,
\]
and says ``hence'' $-u\!\cdot\!K=\Ehat$. The manuscript never defines $T$. A global source
search finds no other occurrence that supplies the missing definition. This is the first line
of the central Vaidya invariant subsection, so a referee is likely to stop there.

If $T$ is an auxiliary tangent, define its components and role before the identity and show the
implication. If the controlled constraint is already $-u_v=\Ehat$ with $K=\partial_v$, the
auxiliary identity may be unnecessary and can be deleted. This is a clarity repair, not an
invitation to change equation (12).

\subsection{Unannounced switch from $\Ehat$ to $E$}

\must. The control charge is $\Ehat$ through equations (1), (7)--(18) and Table 1. Section 4.4
then defines the sextic with $DE=(E^2-1)r+2m$, and Appendices A--B continue with $E$. No
sentence states $E\equiv\Ehat$. Figure titles also use $E$.

Define the identification at the first switch or use $\Ehat$ throughout. Consider replacing
the visually ambiguous product-like symbol $DE$ by a named quantity such as $D_E$ or
$\Delta_E$. At minimum, state that it is a single defined polynomial factor rather than the
product of two previously defined variables.

\subsection{Two incompatible definitions of $\varepsilon$}

\must. Table 1 states $\varepsilon=\dot m$. Appendix C and Figure 9 state
$\varepsilon=M\dot m/m$. The latter is dimensionless; the former is the coordinate mass rate.
The expansion $m(v)=m_0(1+\varepsilon s+\cdots)$ then calls $s=v$ but does not say whether $s$
has been scaled by $M$.

Use $\dot m=dm/dv$ for the rate and reserve $\varepsilon=M\dot m/m$ for the dimensionless
bookkeeping parameter. If units $G=c=M=1$ are used in a numerical subsection, say so. If the
slow coordinate is $s=v/M$, print that definition before equation (C.1). This correction is
essential because the central convergence plot is plotted against $\varepsilon$.

\subsection{Collision between the sextic $S$ and the source $\SD$}

\must. Equation (22) defines $\SD(\lambda)=\int_0^\lambda \mathcal D H\,d\lambda'$. The
following paragraph nevertheless writes $S=\int\Theta H_v\,d\lambda'$, and the Figure 9
caption says that the on-shell term drops ``the costate integral $S$''. In the same section,
$S$ is the genus-two spectral sextic under every radical.

Use $\SD$ in the paragraph, caption and any script-facing prose. The visual distinction is
small but the mathematical objects are entirely different. This is precisely the kind of
collision that creates sign or factor misunderstandings during peer review.

\subsection{Path parameter, cost clock and augmented state}

\must. Section 2.1 says the state is only $x=(r,\varphi)$ and the evolution parameter is the
selected clock $s\in\{t,\tau,\eta,v\}$. It also says that all branches share an indicatrix per
unit selector-time and differ in the cost one-form. For a Vaidya proper-time cost, however, the
metric coefficients depend on the advanced time $v$. If $v$ is not the path parameter, it must
be included in the state; if $v$ is the path parameter, $\tau$ is a cost and should not be
listed as the evolution parameter of that Hamiltonian.

Add one protocol paragraph or table that distinguishes:
\begin{enumerate}
  \item the monotone parameter used to describe the path;
  \item the clock one-form being minimised;
  \item the full state, including a time coordinate when the metric is non-autonomous;
  \item the control, already named $\theta$.
\end{enumerate}
The HJB equation currently minimises over $u$, while $u$ already denotes the four-velocity and
later the retarded coordinate. Replace that control variable by $\theta$ or another unused
symbol. This is a formal clarity issue, not merely typography.

\subsection{Incorrect claim that a nome series is printed}

\must. Appendix C says that the genus-one nome series $g^{(1)}$ is ``printed in
Section 4.4''. It is not present there or elsewhere. Section 4.4 also says that the full
Kronecker--Eisenstein evaluation is given in Appendix B, while Appendix B provides the
Weierstrass block form and residue data but not a complete standalone q-series evaluation
recipe.

Choose one honest level of self-containedness:
\begin{itemize}
  \item print the missing series, characteristic, period and normalisation conventions; or
  \item state that the PDF fixes the algebraic reduction and special-function representation,
  while period construction and numerical q-series evaluation require the archived routine.
\end{itemize}
The second option is shorter and already consistent with the final sentence of the same
paragraph. Do not simultaneously claim full reconstruction from the PDF and reliance on an
unprinted routine.

\subsection{The Vaidya algebraic degeneration is still called a separatrix}

\must. The main text and Appendix B correctly say that $J_{\rm deg}$ gives a double root
$r_d<0$ off the physical arc and is not a physical external separatrix. The last paragraph of
Appendix C is nonetheless titled ``Separatrix limit ($J\to J_{\rm deg}$)'' and repeatedly
uses ``separatrix correction'' and ``separatrix form''.

Rename this paragraph ``Algebraic genus-degeneration limit'' and use ``degeneration
correction/form'' for Vaidya. Reserve $J_{\rm sep}$ and the word separatrix for an accessible
dynamical boundary, primarily the rotating companion problem. This completes a correction that
was already requested and only partially implemented.

\subsection{Spatially flat FLRW is not a spatial-curvature effect}

\must. The abstract says the correction ``separates the spatial-curvature effects of the FLRW
base from the radial mass-flow effects''; the introduction similarly contrasts spatial
curvature with radial flow. The chosen FLRW metric is spatially flat, and the paper's actual
logic is that homogeneous time dependence produces freezing without a spatial gradient,
whereas Vaidya introduces radial gradients and mass-flow memory.

Use ``homogeneous expansion'' for FLRW and ``radial spatial-gradient and mass-flow effects'' for
Vaidya. This makes the physical distinction sharper and removes an apparent contradiction from
the abstract.

\subsection{Reference [12] disagrees with its public DOI}

\must. The manuscript and cover letter cite DOI
\href{https://doi.org/10.5281/zenodo.21740000}{10.5281/zenodo.21740000} under the title
``Brachistochrones in conformal Kerr (Thakurta) spacetimes: ergosphere trichotomy,
separatrices, and adiabatic response''. The registered Zenodo/DataCite title is
``Relativistic Brachistochrones in Conformal Kerr Spacetimes: Adiabatic Theory,
Ergosphere Transitions and Hyperelliptic Closed Forms''.

Bibliographic titles must match the retrievable object. Either update the citation and cover
letter to the deposited title, or publish a new companion version with the intended title and
cite that version. If ``Thakurta'' remains in a BibTeX title, protect its capitalization with
braces; the current PDF renders it in lower case.

\subsection{The public snapshots are older than the final build}

\must. The cover letter calls DOI
\href{https://doi.org/10.5281/zenodo.21739998}{10.5281/zenodo.21739998} ``the present paper'',
but the public record contains \path{paper1-10.pdf}, not the 34-page build reviewed here. The
software DOI \href{https://doi.org/10.5281/zenodo.21707378}{10.5281/zenodo.21707378} is active,
findable, v1.0 and MIT licensed, but it was issued before the final production wrappers for
Figures 1 and 9.

Before journal submission, publish immutable final versions or remove claims that older records
are the present files. The safest workflow is: finalise the source; create a new Zenodo version
for Paper I and, if needed, Paper II and the code package; confirm filenames, titles, hashes and
related-identifier links; then update the cover letter and bibliography. Do not mutate a DOI
claim by inference.

\clearpage
\section{Section-by-section editorial assessment}

\subsection{Title and first page}

The title is accurate, searchable and consistent with the Paper-I Zenodo metadata. It is long,
but every phrase describes a real component. A shorter title could improve memorability, yet a
change now would force another synchronisation of the manuscript, cover letter, metadata and
preprint record. My recommendation is to keep the current title unless the author has a strong
strategic preference.

The first page is clean. The lower whitespace is normal for the IOP class and abstract length;
it is not a production defect. Author name, institution and email are present for a
single-anonymous submission. If double-anonymous review is selected, a separate anonymised PDF
must be generated; the present file is not anonymised.

The abstract is comfortably below IOP's 300-word ceiling, contains no citations or equation
numbers, defines FLRW, and reports method, application and conclusion. It correctly says that
global minimisation is conditional. Two phrases need correction: Kodama should be qualified by
``in spherical symmetry'', and ``spatial-curvature effects of the spatially flat FLRW base''
should become a homogeneous-expansion/spatial-gradient contrast.

\subsection{Section 1: Introduction}

The introduction is concise and much better positioned than in the earlier build. It cites the
stationary optical lineage, direct brachistochrone existence/Morse literature and the recent
Fermat-principle survey. It no longer claims to invent relativistic arrival-time variational
theory. The research gap is now defensible: active maintenance of a selected charge in a
genuinely non-autonomous spacetime, a Kodama selector for spherical dynamics, explicit
Hamiltonians and a complete first-order response.

The final map paragraph should foreground Paper I and mention Paper II once. The manuscript
currently contains roughly two dozen visible companion-paper forward references. Their
cumulative effect is to make Paper I look dependent on a second submission even though its
core is self-contained. The introduction is the right place for the formal companion citation;
later repetitions should be pruned.

The phrase contrasting ``time dependence'' with ``spatial curvature'' should be aligned with
the actual result: FLRW supplies homogeneous expansion without spatial gradients; Vaidya adds
radial geometry and mass flow. This correction will also strengthen the cover letter.

\subsection{Section 2.1: Rail constraint and indicatrix}

The sequence from controlled invariant to velocity indicatrix, domain, endpoint protocol and
Pontryagin reduction is logical. The distinction between a strictly convex closed curve and a
convex velocity set is careful. The statement that convexification would admit averaged
off-shell velocities is useful. The fixed launch event, fixed spatial target and free arrival
clock are now described more carefully than in early versions.

The main weakness is parametric precision. ``Per unit selector-time'', ``selected clock'',
``cost one-form'' and ``evolution parameter'' are used close together without a single master
definition. This becomes consequential for the Vaidya $\tau$ branch. Resolve it in one protocol
box rather than distributing clarifications across later appendices.

Table 1 is now compact, but its first caption sentence still calls both $(\Ehat,J)$ ``conserved
control costates'', whereas $\Ehat$ is introduced primarily as an actively maintained rail
invariant/control charge. More importantly, its $\varepsilon=\dot m$ row conflicts with Figure 9
and Appendix C. Correct that row. The final note about conformal normalisations is sufficient;
no further Paper-II notation belongs here.

Figure 1 is now legible and appropriately limited to the two spherical applications. The
caption should end after explaining FLRW and Vaidya. The sentence that the Thakurta--Kerr
indicatrix is ``retained separately'' describes the author's file management, not information
available to a journal reader.

\subsection{Section 2.2: existence, normality and global optimality}

This section should remain substantially unchanged. The compact regular domain, two-sided
coercivity bound, degeneration exclusions and support-function normality argument are visible.
The manuscript does not pretend that a generic PMP extremal is globally minimal. It explains
the role of an HJB verification function and the first Maxwell/conjugate obstruction.

Three editorial refinements remain. First, define the abbreviations PMP and HJB explicitly at
first use. Second, in Theorem I.2 the symbol $v$ denotes a velocity argument of $F(x,v)$ while
$v$ elsewhere is Vaidya advanced time; use a neutral tangent variable. Third, equation (6)
minimises over $u$, colliding with four-velocity and retarded time; use the already-defined
control $\theta$.

The rotating ergosphere paragraph is useful as a domain caveat but is longer than Paper I needs.
One sentence can state that the timelike-selector hypothesis excludes regions treated in the
companion paper. The repeated details of analytic continuation and alternative congruences are
not necessary to establish the Vaidya theorem.

\subsection{Section 2.3: selector hierarchy}

Recovery of Perlick's fixed-energy stationary optics is an important consistency check. The
hierarchy is stated honestly: Killing where available, conformal Killing in conformally
stationary settings, Kodama in spherical dynamics, and a declared convention otherwise. Keep
that final limitation visible so the abstract does not imply a universal Kodama construction.

The derivation of the Randers data is appropriately technical for CQG. The conformal-Killing
paragraph should serve FLRW directly; the Thakurta--Kerr sentence may be reduced to a citation.

\subsection{Section 3: FLRW}

FLRW functions effectively as the degenerate base: homogeneous expansion changes the local
speed and produces freezing but does not create spatial branch splitting. Equation (10) states
the shared comoving straight line, and Figure 2 is correctly described as a local numerical
check rather than the proof. Preserve the analytical/numerical distinction.

The section is compact enough. The only recommended clarification is to state the endpoint and
regularity assumptions immediately around the three-way $\arg\min$ equality, so it is not read
as a claim about every global topology or target set.

\subsection{Section 4.1: Kodama rail and Hamiltonians}

This is the conceptual heart of Paper I. The unnormalised Kodama choice, controlled-versus-Noether
distinction, branch Hamiltonians and costate/mechanical-angular-momentum map are all valuable.
The exact relation (16) makes the control meaning of $J$ concrete.

Fix the undefined $T$ before submission. Also define prime notation at equation (13): the
reader should not have to infer which clock parametrises $r'$ and $\varphi'$. The paragraph on
on-shell versus off-shell departure is an effective bridge to Theorem I.5, but the surrounding
Thakurta--Kerr implementation details can be reduced.

Figure 3 is a strong numerical sanity check. Its residual and meaning are stated in the caption.
Figure 4 contains useful trajectory and costate information, but the parameter values are mainly
inside the plot. Move the essential values to the caption for accessibility and indexing.

\subsection{Sections 4.2--4.3: phenomenology and frozen-clock comparison}

Penetration, arrival-epoch sensitivity and regular integration through periapsis give the paper
physical texture. The caveat that negative-rate ingoing Vaidya is formal is visible in prose,
captions and legends. The absence-of-inversion result is correctly limited to a frozen-clock
comparison, not presented as a solved outgoing evaporation problem.

The claims should remain tiered. ``Accretion opens a finite window'' and the timing trends appear
to be numerical/model-family results; captions or prose should state the parameter regime and
avoid sounding like parameter-free theorems. Spell out ``boundary-value problem (BVP)'' before
the first acronym.

Figures 5--8 are readable but visually heterogeneous. Most use colour as the principal branch
identifier and rely on small internal titles for parameters. Distinct line styles/markers and
self-contained captions would improve print robustness without changing content.

\subsection{Section 4.4: complete first-order response}

The section now has the correct narrative pivot: the mass derivative of the frozen family is
only the on-shell component; Theorem I.5 supplies the costate-driven off-shell component; the
Vaidya source reduces to a boundary form; and Figure 9 shows that the full response closes the
true-flow residual quadratically. That logical chain should be the main section's dominant
message.

The section remains the largest desk-review risk because it interleaves four layers:
\begin{enumerate}
  \item physical clock comparison and horizon behaviour;
  \item generic first-order canonical perturbation theory;
  \item genus-two Abelian/iterated-integral representation;
  \item the negative-radius algebraic degeneration and Paper-II analogies.
\end{enumerate}
CQG can accommodate advanced mathematics, but the broad reader should be able to leave the
main text knowing the physical result before encountering divisors, local uniformisers and
nome-series implementation. A modest compression would help: keep equations (19)--(24),
Theorem I.5, the boundary source, the representation hierarchy and Figure 9; move repeated
degeneration and companion implementation discussion to the appendices.

Do not weaken the explicit conjectural labels. The paper correctly admits that weight-one
reducibility, a canonical single-valued completion and the minimal basis dimension remain open.
Repair, however, the false statement that $g^{(1)}$ is printed and the collision between $S$ and
$\SD$.

\subsection{Section 5: Conclusions}

The conclusion is concise, uninterrupted and logically tiered. ``Proved'' includes the regular
timelike-selector qualification. ``Conditional and numerical'' contains the ingoing/frozen
limits. ``Outlook'' identifies the full outgoing problem and companion rotation. This is the
right ending architecture.

Replace the spatial-curvature wording if it reappears during copy-editing. Keep the last
companion-paper mention brief. The strongest final statement belongs to Paper I: non-stationary
mass flow changes the controlled arrival problem through a genuine off-shell memory term that
can be computed and tested.

\subsection{Appendix A: reproducibility}

Appendix A is now publication-ready in layout. Table A1 distinguishes exact algebraic reduction
from displayed decimals. Table A2 maps results to scripts and includes the off-shell hash without
collision. Figures A1 and A2 are explicitly local/stationary checks rather than universal
optimality proofs.

The remaining task is archive alignment. If the final source is released as v1.1 or v2, update
Table A2 and the data citation only where the actual script or hash changes. The frozen numerical
script used for Figure 9 should remain identifiable by its current hash.

\subsection{Appendix B: algebraic genus degeneration}

Appendix B is mathematically substantial and now states that $r_d<0$ is not a physical accessible
separatrix. Its editorial role is an analytic degeneration check and an explicit genus-one limit,
not Vaidya phenomenology. Begin and end the appendix with that role. Avoid importing the rotating
separatrix vocabulary.

The residue formula (B.2) now fits the page. The main concern is whether the q-series evaluation
is actually specified or only named. Align the text with what a reader can reconstruct.

\subsection{Appendix C: extended-Hamiltonian derivation}

Appendix C is central and should remain. It supplies the derivation promised by Theorem I.5,
states the turning-point non-uniformity and distinguishes metric-modulus drift from conformal
momentum dilation. The Vaidya boundary reduction is an attractive structural result.

The final ``Separatrix limit'' heading is the outstanding notation regression. Rename it. The
path orientation and dimensionless slow variable also need one clean definition near equation
(C.1), not piecemeal sign explanations several pages later. The ``What is reconstructible''
paragraph should become the definitive and internally consistent statement of closedness.

\subsection{Back matter}

The data statement names an active DOI, version and licence and uses the DOI as the primary
persistent identifier. The acknowledgement identifies the AI models and versions, their uses,
critical review and author responsibility. This is substantively aligned with IOP's current
generative-AI policy. Funding, competing interests and ORCID are present.

Before submission, verify that the University of Pavia affiliation used in the manuscript,
cover letter, ORCID and Zenodo records is accurate and authorised. The code-release DataCite
metadata currently carries a different affiliation string from the preprint records; metadata
consistency is preferable even when the underlying employment/academic situation is legitimate.

\clearpage
\section{Literature positioning and bibliography}

\subsection{Positioning is now adequate}

The introduction now contains the essential lineage: Perlick's fixed-energy optics; Kovner and
Randers/Zermelo/Finsler geometry; the Giannoni--Piccione--Verderesi/Tausk work on relativistic
brachistochrones, existence and Morse theory; and the recent Caponio--Javaloyes survey. This
repairs the previous vulnerability that the paper appeared to claim the whole arrival-time
problem as new.

The novelty should remain phrased narrowly:
\begin{quote}
The new contribution is the active controlled maintenance of a selected rail charge in a
genuinely non-autonomous spacetime, its canonical Kodama realisation in spherical symmetry,
the resulting free-arrival Hamiltonian system in FLRW and Vaidya, and a complete first-order
Vaidya response including the off-shell costate term.
\end{quote}

The special-function literature is sufficient for an editorial screen: Baker, Fay,
Bloch/Zagier, Beilinson--Levin, Brown--Levin and recent higher-genus iterated-integral references
are present. A mathematical referee may request a sharper definition of ``genus-two
dilogarithm'' or additional priority citations. The manuscript already calls that terminology
descriptive and marks irreducibility claims conjectural, which reduces editorial risk.

\subsection{Citation audit}

The final bibliography contains 43 entries and every entry is cited. The former shared-database
inflation has not returned. The visible References heading appears once. Cross-references resolve.
This is a strong improvement.

Several entries remain weaker than the surrounding list:
\begin{itemize}
  \item Reference [12] has the wrong public title and lower-cases ``Thakurta''.
  \item Reference [15] is rendered only as authors, year, \emph{Mem. Amer. Math. Soc.} 300 and
  an arXiv number. The source BibTeX contains the title, issue 1501 and DOI
  \href{https://doi.org/10.1090/memo/1501}{10.1090/memo/1501}; make the visible entry
  independently locatable.
  \item Reference [31] is rendered as ``arXiv preprint'' without its title. It should visibly
  say Brown and Levin, ``Multiple elliptic polylogarithms'', arXiv:1110.6917.
  \item Reference [32] should match the official Zenodo software-record title, or the record
  metadata should be updated. Keep version v1.0 and exact commit only if that is the archive
  actually supporting the submitted paper.
  \item Reference [39] is a patched software tree. The archived snapshot is more important than
  the live GitHub URL; typeset \texttt{is\_*} unambiguously and cite the precise archive.
\end{itemize}

IOP's style guide says article titles are optional for most journals, but a reference should
give enough information to locate the work and include a permanent identifier when available.
The old \texttt{iopart-num} bibliography style suppresses many DOI fields already present in the
source. Production may normalise references later, but the author should perform a verified DOI
pass before submission. Do not infer DOIs from titles. Prioritise the direct brachistochrone,
Kodama/Vaidya, dynamical-horizon and software references.

\subsection{Public-record audit}

The three Zenodo identifiers checked are active and findable:
\begin{itemize}
  \item Paper I preprint: \href{https://doi.org/10.5281/zenodo.21739998}{10.5281/zenodo.21739998},
  currently containing \path{paper1-10.pdf};
  \item Paper II preprint: \href{https://doi.org/10.5281/zenodo.21740000}{10.5281/zenodo.21740000},
  whose registered title differs from reference [12];
  \item software v1.0: \href{https://doi.org/10.5281/zenodo.21707378}{10.5281/zenodo.21707378},
  a 30 July 2026 MIT-licensed archive linked to GitHub tag v1.0.
\end{itemize}

These are real persistent objects, not invented citations. The problem is version alignment.
The final journal package should not call an older PDF ``the present paper''. Create a new
Zenodo version after the last edit and use the version DOI in the cover letter if exact
immutability matters; retain the concept DOI where a moving all-versions citation is intended.

\clearpage
\section{Figures, tables and page design}

\subsection{Global visual assessment}

The current PDF is visually sound. No page contains clipped text, overlapping table cells,
broken glyphs or an isolated float artefact. The previous near-empty pages have disappeared.
The 34-page count is within the production target. Headers, page numbers, section hierarchy and
equation placement are consistent.

Figure 1 is now a successful journal figure: two relevant panels, page-width vector output,
legible labels and differentiated line styles. Figure 9 is also successful: vector, grayscale
safe, compact, and quantitatively documented. The old equation-number mismatch in its plot has
been corrected to equation (22).

Figures 2--8 and A1--A2 remain a mixture of raster outputs and research-diagnostic styles. They
are readable at 100 percent, but they rely heavily on colour, internal titles and legends whose
font size varies. CQG accepts colour, yet a consistent visual language would make the paper look
less like a compilation of separate computational notebooks.

\subsection{Figure-by-figure recommendations}

\begin{longtable}{@{}p{0.12\textwidth}p{0.27\textwidth}p{0.53\textwidth}@{}}
\toprule
Figure & Status & Recommendation \\
\midrule
\endhead
1 & Strong & Remove the file-management sentence about TK being retained separately. Add explicit
panel labels (a), (b) if the plot itself is revised again. \\
2 & Good & Keep the analytical/local-check distinction. Ensure panel labels and final fonts survive
production scaling. \\
3 & Good diagnostic & Preserve the residual threshold. A shorter in-plot title and vector export
would improve consistency. \\
4 & Useful & Put essential $(E,J,m',v_1,r_1)$ values in the caption; use line styles as well as
colour. \\
5 & Useful & State the parameter regime and whether curves are fits, exact frozen laws or numerical
non-autonomous results. \\
6 & Qualified result & Keep ``frozen-clock'' and ``not outgoing BVP'' in the caption, but avoid
making a negative result occupy more visual emphasis than the main Vaidya response. \\
7 & Strong physical illustration & Preserve the formal-negative-rate label. Differentiate all three
orbits by pattern/marker and give units/conventions. \\
8 & Strong timing illustration & Carry arrival and mass-law parameters in the caption; markers should
remain distinguishable in grayscale. \\
9 & Publication-ready & Preserve the current PDF, caption data, exact slopes and original numerical
script hash. \\
A1 & Good check & Vector export desirable; current caption correctly says stationary-limit check. \\
A2 & Good local check & Preserve ``within the tested family'' and avoid any return to global
``minimises'' language. \\
\bottomrule
\end{longtable}

\subsection{Tables}

Table 1 is legible and useful, but the $\varepsilon$ row must be corrected. Table A1 is clear.
Table A2 now wraps correctly; the script name, slope and SHA-256 line no longer collide. Keep
monospaced paths and do not widen the table again.

\subsection{PDF production}

The PDF metadata contains the exact title, author, CQG subject and a useful keyword set. All
fonts are embedded. There are a few benign underfull-box warnings from bibliography URLs and
normal paragraph justification, but no overfull content. The final build should repeat these
checks after notation and citation repairs.

\clearpage
\section{Notation and terminology audit}

\begin{longtable}{@{}p{0.18\textwidth}p{0.31\textwidth}p{0.43\textwidth}@{}}
\toprule
Item & Current use & Required or recommended action \\
\midrule
\endhead
$\Ehat$ versus $E$ & Controlled charge in early sections; bare $E$ in the sextic and appendices. &
Define $E\equiv\Ehat$ once or use one symbol. \must \\
$\varepsilon$ & Both $\dot m$ and $M\dot m/m$. & Reserve $\dot m$ for the rate and $\varepsilon$
for the dimensionless parameter; scale the slow clock explicitly. \must \\
$S$ versus $\SD$ & Spectral sextic and off-shell source. & Use $\SD$ for the source everywhere,
including Figure 9. \must \\
$T$ & Appears in $g(T,T)$ only. & Define or delete the auxiliary identity. \must \\
$u$ & Four-velocity, HJB control and retarded time. & Keep $u^a$ for four-velocity, $u$ for retarded
time only in the outgoing subsection, and use $\theta$ for control. \must \\
$v$ & Advanced time and tangent variable in $F(x,v)$. & Rename the generic tangent variable. \should \\
$s$ & Selected evolution clock and slow parameter. & Distinguish path parameter from dimensionless
slow coordinate and optimised clock. \must \\
$\lambda$ & Elapsed clock from the free endpoint, oriented inward. & Define orientation at first use
in the main theorem and repeat compactly in Appendix C. \should \\
$J_{\rm pen}$ & Physical/dynamical penetration threshold. & Current distinction is good; preserve. \\
$J_{\rm deg}$ & Algebraic double-root value at $r_d<0$. & Call it genus degeneration, not Vaidya
separatrix. \must \\
$J_{\rm sep}$ & Potential accessible spectral separatrix, mainly Paper II. & Do not introduce in
Paper I unless needed; reserve the word. \\
$DE$ & Polynomial factor $(E^2-1)r+2m$. & Prefer $D_E$ or $\Delta_E$, or explicitly state that
$DE$ is one named factor. \should \\
$F$ & Optical cost, orbit integrand and later rational function. & Reuse is survivable but a subscript
would reduce cognitive load in Appendices B--C. \should \\
PMP/HJB/BVP & Used as abbreviations without a parenthetical definition. & Define each at first use;
IOP asks acronyms to be defined in abstract and main text. \should \\
``closed form'' & Sometimes means an evaluable representation requiring period construction/code. &
Use ``closed special-function representation'' where appropriate and state the numerical layer honestly. \\
``separatrix'' & Sometimes physical, sometimes algebraic degeneration. & Reserve for a dynamical
boundary; use ``genus-degeneration locus'' for Vaidya $r_d<0$. \must \\
\bottomrule
\end{longtable}

Hyphenation and spelling are otherwise consistent: ``non-stationary'', ``controlled-rail'',
``first-order'', British ``minimisation'', and ``proper time'' are stable. Equation punctuation
is generally good. Parenthetical density, rather than grammar, is the remaining prose issue.

\clearpage
\section{Scientific narrative and claim hierarchy}

The paper now has a viable six-question narrative:

\begin{enumerate}
  \item \textbf{What is controlled?} A fixed observer/selector-relative rail charge restricts
  the admissible timelike velocities to a compact indicatrix on the regular domain.
  \item \textbf{Why is the optimisation legitimate?} Pontryagin gives the extremal system;
  existence and normality hold under explicit compactness/coercivity hypotheses; global
  optimality requires a separate HJB certificate.
  \item \textbf{Which geometric field supplies the charge?} Killing in stationary geometry,
  conformal Killing in the FLRW base, and Kodama in spherical dynamics.
  \item \textbf{What does time dependence do?} FLRW gives homogeneous freezing without branch
  splitting; Vaidya adds radial wind, costate memory, penetration and timing effects.
  \item \textbf{Why is the frozen derivative insufficient?} The non-autonomous orbit leaves the
  instantaneous shell through an interior costate; Theorem I.5 supplies that off-shell term.
  \item \textbf{Can it be evaluated and checked?} The response has a finite Abelian/iterated
  representation and the true-flow residual changes from first- to second-order scaling.
\end{enumerate}

This order should dominate the abstract, introduction, section headings and conclusion. The
negative-radius degeneration and companion Kerr dilation are supporting structures, not equal
headline results for Paper I.

The claim-status vocabulary is mostly disciplined. The following wording should remain stable:
\begin{itemize}
  \item ``PMP extremal'' or ``candidate minimizer'' unless a verification function or global
  comparison is actually supplied;
  \item ``local perturbation check'' for Figures 2 and A2;
  \item ``frozen-clock comparison'' for the advanced/retarded and no-inversion result;
  \item ``formal negative-rate continuation'' for ingoing $m'<0$;
  \item ``closed special-function representation'' when numerical periods are still needed;
  \item ``conjectural'' for irreducibility/minimal-basis/single-valued-completion statements;
  \item ``open'' for the full outgoing-Vaidya BVP and generic global optimality.
\end{itemize}

The word ``universal'' appears in several contexts. It is justified for the Vaidya boundary
source only within the stated self-similar Hamiltonian class and canonical conventions. Avoid
allowing the cover letter or abstract to broaden it to arbitrary non-stationary metrics.

\clearpage
\section{Reproducibility, data, software and AI policy}

CQG has adopted IOP's research-data policy. IOP requires a data-availability statement and asks
authors to cite public software/data with persistent identifiers. The manuscript complies in
form: it has a data statement, an active software DOI, version, licence, GitHub development link,
script map and software references.

The software DOI metadata confirms an MIT-licensed v1.0 release issued on 30 July 2026 and linked
to GitHub tag v1.0. Its description includes the controlled-rail Hamiltonians, first-order
on/off-shell corrections, regression runner and Sage requirements. This supports the scientific
claims. The final manuscript should, however, cite the exact version that includes the submitted
source and final production scripts, or explicitly distinguish scientific computation from
later editorial plotting.

The current AI acknowledgement is acceptable in substance under IOP's August 2026 ethical
guidance. It identifies OpenAI GPT-5.6 (Sol) and Anthropic Claude (Opus 4.8), says they assisted
with symbolic-algebra checks, numerical-code development and language/structural editing, and
states independent verification and author responsibility. Keep those four components. Do not
list an AI system as an author. Do not let AI-origin bibliography entries remain unverified.

Because the final Figure 1 and Figure 9 production wrappers were created after software v1.0,
the acknowledgement is not the issue; archive completeness is. A final release should include:
\begin{itemize}
  \item the submitted \LaTeX{} source and bibliography;
  \item the exact figure-generation wrappers and frozen numerical source;
  \item environment/lock information;
  \item the script hash quoted in Table A2;
  \item a short README mapping paper figures/equations to scripts;
  \item the final 34-page PDF or its precise preprint version.
\end{itemize}

The cover letter is strong in scope and contribution, but its DOI paragraph must be synchronised.
The title it gives Paper II does not match the public DOI, and ``the present paper is DOI ...''
points to an older PDF. Correct these before ScholarOne submission. The request for coordinated
handling is reasonable, provided both papers are independently intelligible and their public
records match their submitted titles.

\clearpage
\section{Major and minor action checklists}

\subsection{Mandatory pre-submission checklist}

\begin{enumerate}[label=\textbf{M\arabic*.}]
  \item Define or remove $T$ in the first line of section 4.1.
  \item Define $E\equiv\Ehat$ or standardise the energy symbol.
  \item Standardise $\dot m$ and $\varepsilon=M\dot m/m$; define the dimensionless slow clock.
  \item Replace the off-shell source $S$ by $\SD$ in prose and Figure 9.
  \item Clarify path parameter, cost clock, non-autonomous state augmentation and control symbol.
  \item Remove or repair the false claim that $g^{(1)}$ is printed; calibrate text-only
  reconstructibility.
  \item Rename the Appendix-C Vaidya ``Separatrix limit'' as algebraic genus degeneration.
  \item Replace the spatial-curvature/FLRW wording in abstract and introduction.
  \item Make reference [12] and the cover letter match the public Paper-II DOI title.
  \item Publish or cite final matching Zenodo versions; do not call \path{paper1-10.pdf} the
  present 34-page paper.
\end{enumerate}

\subsection{Strongly recommended editorial checklist}

\begin{enumerate}[label=\textbf{S\arabic*.}]
  \item Reduce companion-Paper-II forward references to approximately five strategically placed
  mentions.
  \item Remove the ``retained separately'' sentence from Figure 1's caption.
  \item Define PMP, HJB and BVP at first use.
  \item Rename the generic tangent argument of $F(x,v)$ to avoid collision with Vaidya time.
  \item Move essential parameters from Figures 4--8 plot titles into captions.
  \item Add line styles/markers and vector exports to remaining multi-curve figures where practical.
  \item Make references [15], [31], [32] and [39] independently locatable in the rendered list.
  \item Run a verified DOI pass without guessing identifiers.
  \item Compress repeated special-function and Paper-II implementation prose in section 4.4 and
  Appendices B--C.
  \item Verify consistent author affiliation across manuscript, ORCID, cover letter and Zenodo.
\end{enumerate}

\subsection{Final production checklist}

\begin{enumerate}[label=\textbf{P\arabic*.}]
  \item Full \texttt{pdflatex/bibtex} build with zero undefined references and zero overfull boxes.
  \item Search the PDF and build files for ``??'', stale equation numbers and duplicate labels.
  \item Confirm Theorem I.5 remains in the main text and equation (22) remains the Figure-9 target.
  \item Confirm Table A2 and equation (B.2) remain within margins.
  \item Re-render all pages and inspect figures, bibliography and page transitions.
  \item Confirm PDF title, author, subject, keywords and embedded fonts.
  \item Confirm the final DOI links resolve to the exact titles and versions stated.
  \item Prepare the chosen single- or double-anonymous package consistently.
\end{enumerate}

\clearpage
\section{What should not be reopened}

The final cleanup should be narrow. Do not destabilise the manuscript by rederiving settled
content without a concrete mathematical reason. Preserve:

\begin{itemize}
  \item the controlled invariant $-u\!\cdot\!W=\Ehat$ as the paper's central definition;
  \item the selector hierarchy with its spherical-symmetry limitation;
  \item the operational rail-force and control-power discussion;
  \item the compact strictly convex indicatrix assumptions;
  \item the fixed launch/fixed spatial target/free arrival protocol;
  \item the distinction between endpoint transversality and an interior first integral;
  \item the costate/mechanical-angular-momentum distinction;
  \item the conditional HJB and cut/Maxwell caveats;
  \item FLRW as the homogeneous branch-degenerate base;
  \item the unnormalised Kodama selector for ingoing Vaidya;
  \item the formal status of negative-rate ingoing continuation;
  \item Theorem I.5 in the main text and its derivation in Appendix C;
  \item the turning-point non-uniformity statement;
  \item the Vaidya boundary source $\SD=[rp_r]-\lambda$ after notation cleanup;
  \item Figure 9 and its exact regression values;
  \item explicit conjectural/open labels for the special-function hierarchy;
  \item the ``Proved / Conditional and numerical / Outlook'' conclusion;
  \item the data DOI, script map and transparent AI disclosure, updated only for exact versioning.
\end{itemize}

\clearpage
\section{Recommended final revision sequence}

The safest sequence is short and mechanical:

\begin{enumerate}
  \item \textbf{Repair definitions first.} Fix $T$, $E/\Ehat$, $\varepsilon$, $S/\SD$,
  path/cost/state conventions and control notation. Recompile and check equations (12), (22),
  (C.1)--(C.6).
  \item \textbf{Repair the special-function claims.} Remove the nonexistent printed $g^{(1)}$
  claim, choose the honest reconstructibility wording, and rename the algebraic degeneration.
  \item \textbf{Repair the public records.} Decide final titles, create new Zenodo versions as
  needed, then update reference [12], reference [32] and the cover letter from verified metadata.
  \item \textbf{Perform the bibliography pass.} Improve [15], [31] and [39], add only verified
  permanent identifiers and keep all 43 entries cited.
  \item \textbf{Trim Paper-II leakage.} Remove redundant forward comparisons without touching
  Paper-I derivations.
  \item \textbf{Polish captions and plots.} At minimum make parameters self-contained; vectorise
  only if this does not risk numerical content.
  \item \textbf{Build and render.} Repeat the clean 34--35-page production checks and inspect every
  page.
  \item \textbf{Submit.} Upload the complete PDF, relevant supplementary/code information and the
  corrected cover letter; select the intended peer-review anonymity model.
\end{enumerate}

This sequence should take a focused editorial pass, not a new research cycle. The manuscript
should not be delayed for the full outgoing-Vaidya problem, a proof of higher-genus
irreducibility, or a global HJB solution, because it already labels those as open.

\clearpage
\section{Final recommendation}

The present build is substantially stronger than the versions reviewed on 3 August 2026. It is
within CQG's scope, contains a plausible significant advance, has a defensible novelty position,
and presents its principal off-shell correction in the main text with a convincing numerical
order test. The visual production is clean and the reproducibility posture is strong.

It is not yet the version I would upload. The undefined $T$, unannounced $E/\Ehat$ switch,
inconsistent $\varepsilon$, source/sextic symbol collision, ambiguous non-autonomous state
formulation, false ``printed series'' claim, residual separatrix terminology and DOI/version
mismatches are all concrete. Leaving them in would create avoidable doubt precisely where the
paper asks a referee to trust a difficult formal construction.

After those corrections, my simulated recommendation is:

\begin{center}
\colorbox{LightBlue}{\parbox{0.90\textwidth}{
\textbf{Send for external peer review as a CQG Research Paper.}
The paper is not editorially acceptance-ready---no pre-submission assessment can replace
specialist review---but it should be desk-review-ready after the short mandatory pass. Its
remaining risk will then be scientific rather than presentational: whether referees accept the
optimal-control foundations, the first-order canonical signs and the claimed special-function
closure.
}}
\end{center}

In direct answer to ``is it ready?'': \textbf{almost, but not today}. Correct the ten mandatory
items, synchronise the Zenodo records, and submit the resulting build. Further broad rewriting is
not justified.

\vspace{0.45em}
{\large\bfseries Sources consulted\par}
\addcontentsline{toc}{section}{Sources consulted}

\begingroup
\small
\raggedright
\setlist[itemize]{itemsep=0.12em,topsep=0.2em}
\begin{itemize}
  \item Current manuscript: \path{paper1/paper1.pdf}, 34 pages, built 4 August 2026.
  \item Supplied brief: \path{Instructions_for_Codex_CQG_Editorial_Review.pdf}, 3 August 2026.
  \item CQG scope, article types, peer review and data-policy page:\\
  \url{https://publishingsupport.iopscience.iop.org/journals/classical-and-quantum-gravity/about-classical-quantum-gravity/}
  \item IOP article structure guidance:\\
  \url{https://publishingsupport.iopscience.iop.org/questions/structure-and-format-of-your-journal-article/}
  \item IOP style guide for journal articles:\\
  \url{https://publishingsupport.iopscience.iop.org/questions/style-guide-journal-articles/}
  \item IOP submission checklist:\\
  \url{https://publishingsupport.iopscience.iop.org/submission-checklist/}
  \item IOP research-data policy:\\
  \url{https://publishingsupport.iopscience.iop.org/iop-publishing-data-availability-policy/}
  \item IOP ethical policy, including generative-AI guidance:\\
  \url{https://publishingsupport.iopscience.iop.org/ethical-policy-journals/}
  \item Caponio and Javaloyes, arXiv:2605.01532:\\
  \url{https://arxiv.org/abs/2605.01532}
  \item Brown and Levin, arXiv:1110.6917:\\
  \url{https://arxiv.org/abs/1110.6917}
  \item DataCite/Zenodo records for DOI 10.5281/zenodo.21739998,
  10.5281/zenodo.21740000 and 10.5281/zenodo.21707378, accessed 4 August 2026.
\end{itemize}
\endgroup

\end{document}
